% Infografía de la dualidad en una cara (Víctor, 15/11/2026): amplía el esquema de
% dualidad (figura esquema-dualidad) con dos paneles en 3D, en la perspectiva del
% gráfico pmu-r3 del 3.A.8 (acimut 30°, elevación 28°):
%  - primal: superficie de utilidad U = (x1 x2)^(1/2), recta de balance en el suelo y
%    plano vertical de la restricción de gasto; la utilidad recorrida sobre la recta
%    alcanza su máximo en x* (U* = V(p,W));
%  - dual: el gasto p·h como plano inclinado, la curva de indiferencia exigida U-barra
%    en el suelo y su «cortina» hasta el plano; el gasto recorrido sobre la curva
%    alcanza su mínimo en h* (E* = E(p,U)).
% Debajo, el esquema ampliado: problemas, demandas, funciones valor, Roy, Shephard,
% identidades de dualidad, Slutsky, multiplicadores, propiedades y aplicaciones.
% Datos: p1 = p2 = 1, W = 3,5; x* = h* = (1,75; 1,75); U* = U-barra = 1,75; E* = 3,5.
% Primal (marshalliana) en morado y dual (hicksiana) en verde, como el esquema.
% Autor: Víctor Gutiérrez Marcos
\documentclass[tikz,border=4pt]{standalone}
% ==========================================================================
% tikz-tcee.tex — Estilos comunes de los gráficos TikZ del temario
% Autor: Víctor Gutiérrez Marcos
%
% Cada gráfico es un fichero standalone en fuentes/<tema>/graficos/:
%
%   \documentclass[tikz,border=4pt]{standalone}
%   % ==========================================================================
% tikz-tcee.tex — Estilos comunes de los gráficos TikZ del temario
% Autor: Víctor Gutiérrez Marcos
%
% Cada gráfico es un fichero standalone en fuentes/<tema>/graficos/:
%
%   \documentclass[tikz,border=4pt]{standalone}
%   % ==========================================================================
% tikz-tcee.tex — Estilos comunes de los gráficos TikZ del temario
% Autor: Víctor Gutiérrez Marcos
%
% Cada gráfico es un fichero standalone en fuentes/<tema>/graficos/:
%
%   \documentclass[tikz,border=4pt]{standalone}
%   \input{../../../latex/tikz-tcee}
%   \begin{document}
%   \begin{tikzpicture}[tcee]
%     \ejes{Q}{P}{6}{5}
%     \draw[curva1] (0.5,4.5) -- (5,0.8) node[etiqueta,right] {$D$};
%     \capa{2}{\draw[curva2] ...;}   % en el vídeo aparece en el paso 2
%   \end{tikzpicture}
%   \end{document}
%
% Paleta: la de la web (morado, dorado, salvia) más tres colores de apoyo
% distinguibles en blanco y negro por el trazo.
% ==========================================================================
\usepackage[T1]{fontenc}
\usepackage{newpxtext,newpxmath}
\usepackage{xcolor}
\usepackage{tikz,pgfplots}
\pgfplotsset{compat=1.18}
\usetikzlibrary{arrows.meta,calc,positioning,patterns,decorations.pathreplacing,intersections,shapes.misc,backgrounds}

\definecolor{tceemorado}{HTML}{5F2987}
\definecolor{tceedorado}{HTML}{B8860B}
\definecolor{tceeverde}{HTML}{3C7D3A}
\definecolor{tceeazul}{HTML}{1F5F99}
\definecolor{tceerojo}{HTML}{B22222}
\definecolor{tceenaranja}{HTML}{E07B28}
\definecolor{tceegris}{HTML}{767171}
\definecolor{tceesalvia}{HTML}{E2EFD9}

% Capas para el vídeo: \capa{n}{…} dibuja su contenido solo a partir del paso n.
% En el PDF, la web y Word se ve todo; generar-video.py compila el gráfico una
% vez por paso con \def\tceepaso{k} para que cada elemento aparezca al narrarlo.
% Vectores con flecha encima, como en el texto (tcee.sty): \vect{x}, \vect{p}
\providecommand{\vect}[1]{\vec{#1}}
\providecommand{\tceepaso}{1000}
\newcommand{\capa}[2]{\ifnum#1>\tceepaso\relax\else#2\fi}
% Gráficos por pasos del mismo tamaño (propuesta del 3.A.29, 11/11/2026): con
% \def\tceefantasma{1} antes del \input (p. ej., x-paso1.tex:
% «\def\tceepaso{3}\def\tceefantasma{1}\input{x.tex}»), las capas posteriores
% al paso se dibujan invisibles y ocupan su sitio. Sin \tceefantasma (el vídeo),
% no cambia nada.
\ifdefined\tceefantasma
  \renewcommand{\capa}[2]{\ifnum#1>\tceepaso\relax\begin{scope}[opacity=0]#2\end{scope}\else#2\fi}
\fi

\tikzset{
  tcee/.style={line cap=round,line join=round,font=\small,>={Stealth[length=2.2mm]}},
  eje/.style={->,thick,black},
  curva1/.style={very thick,tceemorado},
  curva2/.style={very thick,tceeazul},
  curva3/.style={very thick,tceeverde},
  curva4/.style={very thick,tceedorado},
  curvanueva/.style={very thick,tceerojo},
  desplazada/.style={very thick,dashed},
  guia/.style={thin,dashed,tceegris},
  punto/.style={circle,fill=black,inner sep=1.3pt},
  etiqueta/.style={font=\small,inner sep=2pt},
  flecha/.style={->,thick,tceerojo},
  area/.style={fill=tceemorado,fill opacity=0.18,draw=none},
  area2/.style={fill=tceedorado,fill opacity=0.22,draw=none},
}

% \ejes{etiqueta x}{etiqueta y}{largo x}{alto y}
\newcommand{\ejes}[4]{%
  \draw[eje] (0,0) -- (#3,0) node[below left=2pt and 0pt] {#1};
  \draw[eje] (0,0) -- (0,#4) node[left] {#2};
  \node[below left] at (0,0) {$0$};}
% \proyeccion{(x,y)}{etq x}{etq y}: líneas guía a los ejes con etiquetas
\newcommand{\proyeccion}[3]{%
  \draw[guia] let \p1=#1 in (\x1,0) node[below] {#2} |- (0,\y1) node[left] {#3};}

% --------------------------------------------------------------------------
% Paleta «dictamen» (dic…): gráficos de datos de la Guía del dictamen (1D01) y
% de cualquier gráfico de coyuntura que quiera los mismos colores. Un concepto,
% un color en todos los gráficos (CRITERIOS §7, «tres reglas», 24/10/2026);
% tabla completa en fuentes/1D01/_trabajo/paleta-1D01.md. No cambia nada de
% lo anterior: son colores y estilos nuevos.
% Actividad: el PIB (y sus variantes: PIB por persona, PIB mundial), morado.
\definecolor{dicpib}{HTML}{5F2987}
% Demanda interna, gama cálida: consumo (naranjas: privado oscuro, público
% claro), inversión (verde) y existencias (gris con trama, aparte).
\definecolor{dicconspriv}{HTML}{C76A1F}
\definecolor{dicconspub}{HTML}{F5C690}
\definecolor{dicfbcf}{HTML}{2A6B3C}
\definecolor{dicexistencias}{HTML}{B4B4B4}
% Demanda externa y comercio, gama fría (azul claro, que se separa del morado)
\definecolor{dicexterna}{HTML}{6AA6DA}
% Zonas: zona del euro (azul marino), España (rojo), resto de países (grises)
\definecolor{dicze}{HTML}{17406D}
\definecolor{dicespana}{HTML}{C8382E}
\definecolor{dicotro}{HTML}{8A8A8A}
\definecolor{dicotroclaro}{HTML}{B5B5B5}
\definecolor{dicotromasclaro}{HTML}{DADADA}
% Precios: inflación general (negro, continuo) y subyacente (negro,
% discontinuo); componentes: energía (bermellón), alimentos (amarillo) y los
% de la subyacente en azules (bienes industriales claro, servicios oscuro).
\definecolor{dicinflacion}{HTML}{1A1A1A}
\definecolor{dicenergia}{HTML}{D55E00}
\definecolor{dicalimentos}{HTML}{F2D45C}
\definecolor{dicbine}{HTML}{7FB2E0}
\definecolor{dicservicios}{HTML}{1F4E89}
% Trabajo: empleo (verde azulado), horas (el mismo, claro), tasa de actividad
% (verde azulado oscuro), paro (teja); productividad (dorado de la web).
\definecolor{dicempleo}{HTML}{1B7F6B}
\definecolor{dichoras}{HTML}{86CBB8}
\definecolor{dicactividad}{HTML}{0F5247}
\definecolor{dicparo}{HTML}{A23B2A}
\definecolor{dicproductividad}{HTML}{B8860B}
% Hogares: renta disponible (pizarra) y tasa de ahorro (negro); el consumo es
% dicconspriv.
\definecolor{dicrenta}{HTML}{4A5A6A}
\definecolor{dicahorro}{HTML}{1A1A1A}
% Política monetaria: tipos del BCE en azules (depósito, el de la zona del
% euro; operaciones principales, medio; facilidad marginal, claro); balance:
% préstamos a los bancos (dorado claro), resto del activo (gris claro).
\definecolor{dicdeposito}{HTML}{17406D}
\definecolor{dicopf}{HTML}{4F86C6}
\definecolor{dicfmc}{HTML}{9EC5EA}
\definecolor{dicprestamos}{HTML}{D9A520}
\definecolor{dicresto}{HTML}{CFCFCF}
% Ciclo: series adelantadas, coincidentes (el PIB) y retrasadas
\definecolor{dicadelantada}{HTML}{D55E00}
\definecolor{dicretrasada}{HTML}{6AA6DA}
% Referencias: rejilla, previsión, sombreados
\definecolor{dicreferencia}{HTML}{767171}

\tikzset{
  % barra de un componente: \addplot[dicbarra=dicfbcf]
  dicbarra/.style={fill=#1, draw=#1!65!black, line width=0.3pt},
  % existencias: gris con trama, para que se distinga también en gris
  dicbarratrama/.style={fill=#1, draw=#1!65!black, line width=0.3pt,
    postaction={pattern=north east lines, pattern color=white}},
  % línea de una serie; con halo blanco cuando pasa sobre barras
  diclinea/.style={very thick, #1},
  diclineahalo/.style={very thick, #1, preaction={draw=white, line width=3.4pt, solid}},
  % puntos de una serie (capa aparte, después de la línea)
  dicpuntos/.style={only marks, mark=*, mark size=1.8pt,
    mark options={fill=white, draw=#1, line width=0.9pt, solid}},
  % línea del cero o de un umbral, después de las barras
  diccero/.style={thin, black},
  dicumbral/.style={thin, dashed, dicreferencia},
}
% Capas de los gráficos de datos (CRITERIOS §7, «tres reglas»), de abajo
% arriba: cuadrícula → sombreados (dicfondo) → barras y áreas (dicbarras) → ejes, cero y umbrales
% (axis lines) → líneas (main) → puntos (dicpuntos) → rótulos (axis
% descriptions). No basta el orden del código: pgfplots dibuja cuadrícula,
% ejes y marcas en sus propias capas (y area style pone los ejes encima), y
% dentro de un axis los \draw, \fill y \node no admiten capa (van con los
% datos de main); por eso todo lo que se dibuja dentro del axis es un \addplot.
% Uso: \begin{axis}[…, dicapas] (al final de las opciones);
%   barras y áreas: \addplot[on layer=dicbarras, …];
%   sombreados:     \dicsombra[estilo]{x1}{y1}{x2}{y2} (rectángulo, en dicfondo,
%                   bajo las barras y sobre la cuadrícula);
%   cero y umbrales: \dicraya[estilo]{(x1,y1) (x2,y2)} (en axis lines);
%   líneas: \addplot normal (main); puntos: \addplot[on layer=dicpuntos, …]
%   (las marcas de cualquier serie van solas a dicpuntos);
%   rótulos: \dicrotulo[opciones del nodo]{(x,y)}{texto} (en axis descriptions),
%   siempre sin fondo (nada de fill=white: si chocan con una línea, se desplazan;
%   CRITERIOS §7, regla 4).
% Los \capa{n} del vídeo no cambian: siguen envolviendo instrucciones completas.
\pgfdeclarelayer{dicfondo}
\pgfdeclarelayer{dicbarras}
\pgfdeclarelayer{dicpuntos}
\pgfplotsset{
  % (el segundo argumento reparte cuadrícula, ejes y rótulos como el estándar)
  layers/dic/.define layer set={axis background, axis grid, dicfondo, dicbarras,
    axis ticks, axis lines, axis tick labels, pre main, main, dicpuntos,
    axis descriptions, axis foreground}{/pgfplots/layers/standard},
  % leyenda sin caja blanca: toma el fondo del gráfico, también si el gráfico va
  % dentro de un cuadro gris (CRITERIOS §7, regla 4, 26/10/2026)
  dicapas/.style={axis on top=false, set layers=dic, mark layer=dicpuntos,
    legend style={fill=none}},
  dicsuelto/.style={forget plot, stack plots=false, sharp plot, mark=none},
}
\newcommand{\dicraya}[2][]{%
  \addplot[on layer=axis lines, dicsuelto, diccero, #1] coordinates {#2};}
\newcommand{\dicsombra}[5][]{%
  \addplot[on layer=dicfondo, dicsuelto, draw=none, #1]
    coordinates {(#2,#3) (#4,#3) (#4,#5) (#2,#5)};}
\newcommand{\dicrotulo}[3][]{%
  \addplot[on layer=axis descriptions, dicsuelto, draw=none] coordinates {#2}
    node[pos=0, #1] {#3};}

%   \begin{document}
%   \begin{tikzpicture}[tcee]
%     \ejes{Q}{P}{6}{5}
%     \draw[curva1] (0.5,4.5) -- (5,0.8) node[etiqueta,right] {$D$};
%     \capa{2}{\draw[curva2] ...;}   % en el vídeo aparece en el paso 2
%   \end{tikzpicture}
%   \end{document}
%
% Paleta: la de la web (morado, dorado, salvia) más tres colores de apoyo
% distinguibles en blanco y negro por el trazo.
% ==========================================================================
\usepackage[T1]{fontenc}
\usepackage{newpxtext,newpxmath}
\usepackage{xcolor}
\usepackage{tikz,pgfplots}
\pgfplotsset{compat=1.18}
\usetikzlibrary{arrows.meta,calc,positioning,patterns,decorations.pathreplacing,intersections,shapes.misc,backgrounds}

\definecolor{tceemorado}{HTML}{5F2987}
\definecolor{tceedorado}{HTML}{B8860B}
\definecolor{tceeverde}{HTML}{3C7D3A}
\definecolor{tceeazul}{HTML}{1F5F99}
\definecolor{tceerojo}{HTML}{B22222}
\definecolor{tceenaranja}{HTML}{E07B28}
\definecolor{tceegris}{HTML}{767171}
\definecolor{tceesalvia}{HTML}{E2EFD9}

% Capas para el vídeo: \capa{n}{…} dibuja su contenido solo a partir del paso n.
% En el PDF, la web y Word se ve todo; generar-video.py compila el gráfico una
% vez por paso con \def\tceepaso{k} para que cada elemento aparezca al narrarlo.
% Vectores con flecha encima, como en el texto (tcee.sty): \vect{x}, \vect{p}
\providecommand{\vect}[1]{\vec{#1}}
\providecommand{\tceepaso}{1000}
\newcommand{\capa}[2]{\ifnum#1>\tceepaso\relax\else#2\fi}
% Gráficos por pasos del mismo tamaño (propuesta del 3.A.29, 11/11/2026): con
% \def\tceefantasma{1} antes del \input (p. ej., x-paso1.tex:
% «\def\tceepaso{3}\def\tceefantasma{1}\input{x.tex}»), las capas posteriores
% al paso se dibujan invisibles y ocupan su sitio. Sin \tceefantasma (el vídeo),
% no cambia nada.
\ifdefined\tceefantasma
  \renewcommand{\capa}[2]{\ifnum#1>\tceepaso\relax\begin{scope}[opacity=0]#2\end{scope}\else#2\fi}
\fi

\tikzset{
  tcee/.style={line cap=round,line join=round,font=\small,>={Stealth[length=2.2mm]}},
  eje/.style={->,thick,black},
  curva1/.style={very thick,tceemorado},
  curva2/.style={very thick,tceeazul},
  curva3/.style={very thick,tceeverde},
  curva4/.style={very thick,tceedorado},
  curvanueva/.style={very thick,tceerojo},
  desplazada/.style={very thick,dashed},
  guia/.style={thin,dashed,tceegris},
  punto/.style={circle,fill=black,inner sep=1.3pt},
  etiqueta/.style={font=\small,inner sep=2pt},
  flecha/.style={->,thick,tceerojo},
  area/.style={fill=tceemorado,fill opacity=0.18,draw=none},
  area2/.style={fill=tceedorado,fill opacity=0.22,draw=none},
}

% \ejes{etiqueta x}{etiqueta y}{largo x}{alto y}
\newcommand{\ejes}[4]{%
  \draw[eje] (0,0) -- (#3,0) node[below left=2pt and 0pt] {#1};
  \draw[eje] (0,0) -- (0,#4) node[left] {#2};
  \node[below left] at (0,0) {$0$};}
% \proyeccion{(x,y)}{etq x}{etq y}: líneas guía a los ejes con etiquetas
\newcommand{\proyeccion}[3]{%
  \draw[guia] let \p1=#1 in (\x1,0) node[below] {#2} |- (0,\y1) node[left] {#3};}

% --------------------------------------------------------------------------
% Paleta «dictamen» (dic…): gráficos de datos de la Guía del dictamen (1D01) y
% de cualquier gráfico de coyuntura que quiera los mismos colores. Un concepto,
% un color en todos los gráficos (CRITERIOS §7, «tres reglas», 24/10/2026);
% tabla completa en fuentes/1D01/_trabajo/paleta-1D01.md. No cambia nada de
% lo anterior: son colores y estilos nuevos.
% Actividad: el PIB (y sus variantes: PIB por persona, PIB mundial), morado.
\definecolor{dicpib}{HTML}{5F2987}
% Demanda interna, gama cálida: consumo (naranjas: privado oscuro, público
% claro), inversión (verde) y existencias (gris con trama, aparte).
\definecolor{dicconspriv}{HTML}{C76A1F}
\definecolor{dicconspub}{HTML}{F5C690}
\definecolor{dicfbcf}{HTML}{2A6B3C}
\definecolor{dicexistencias}{HTML}{B4B4B4}
% Demanda externa y comercio, gama fría (azul claro, que se separa del morado)
\definecolor{dicexterna}{HTML}{6AA6DA}
% Zonas: zona del euro (azul marino), España (rojo), resto de países (grises)
\definecolor{dicze}{HTML}{17406D}
\definecolor{dicespana}{HTML}{C8382E}
\definecolor{dicotro}{HTML}{8A8A8A}
\definecolor{dicotroclaro}{HTML}{B5B5B5}
\definecolor{dicotromasclaro}{HTML}{DADADA}
% Precios: inflación general (negro, continuo) y subyacente (negro,
% discontinuo); componentes: energía (bermellón), alimentos (amarillo) y los
% de la subyacente en azules (bienes industriales claro, servicios oscuro).
\definecolor{dicinflacion}{HTML}{1A1A1A}
\definecolor{dicenergia}{HTML}{D55E00}
\definecolor{dicalimentos}{HTML}{F2D45C}
\definecolor{dicbine}{HTML}{7FB2E0}
\definecolor{dicservicios}{HTML}{1F4E89}
% Trabajo: empleo (verde azulado), horas (el mismo, claro), tasa de actividad
% (verde azulado oscuro), paro (teja); productividad (dorado de la web).
\definecolor{dicempleo}{HTML}{1B7F6B}
\definecolor{dichoras}{HTML}{86CBB8}
\definecolor{dicactividad}{HTML}{0F5247}
\definecolor{dicparo}{HTML}{A23B2A}
\definecolor{dicproductividad}{HTML}{B8860B}
% Hogares: renta disponible (pizarra) y tasa de ahorro (negro); el consumo es
% dicconspriv.
\definecolor{dicrenta}{HTML}{4A5A6A}
\definecolor{dicahorro}{HTML}{1A1A1A}
% Política monetaria: tipos del BCE en azules (depósito, el de la zona del
% euro; operaciones principales, medio; facilidad marginal, claro); balance:
% préstamos a los bancos (dorado claro), resto del activo (gris claro).
\definecolor{dicdeposito}{HTML}{17406D}
\definecolor{dicopf}{HTML}{4F86C6}
\definecolor{dicfmc}{HTML}{9EC5EA}
\definecolor{dicprestamos}{HTML}{D9A520}
\definecolor{dicresto}{HTML}{CFCFCF}
% Ciclo: series adelantadas, coincidentes (el PIB) y retrasadas
\definecolor{dicadelantada}{HTML}{D55E00}
\definecolor{dicretrasada}{HTML}{6AA6DA}
% Referencias: rejilla, previsión, sombreados
\definecolor{dicreferencia}{HTML}{767171}

\tikzset{
  % barra de un componente: \addplot[dicbarra=dicfbcf]
  dicbarra/.style={fill=#1, draw=#1!65!black, line width=0.3pt},
  % existencias: gris con trama, para que se distinga también en gris
  dicbarratrama/.style={fill=#1, draw=#1!65!black, line width=0.3pt,
    postaction={pattern=north east lines, pattern color=white}},
  % línea de una serie; con halo blanco cuando pasa sobre barras
  diclinea/.style={very thick, #1},
  diclineahalo/.style={very thick, #1, preaction={draw=white, line width=3.4pt, solid}},
  % puntos de una serie (capa aparte, después de la línea)
  dicpuntos/.style={only marks, mark=*, mark size=1.8pt,
    mark options={fill=white, draw=#1, line width=0.9pt, solid}},
  % línea del cero o de un umbral, después de las barras
  diccero/.style={thin, black},
  dicumbral/.style={thin, dashed, dicreferencia},
}
% Capas de los gráficos de datos (CRITERIOS §7, «tres reglas»), de abajo
% arriba: cuadrícula → sombreados (dicfondo) → barras y áreas (dicbarras) → ejes, cero y umbrales
% (axis lines) → líneas (main) → puntos (dicpuntos) → rótulos (axis
% descriptions). No basta el orden del código: pgfplots dibuja cuadrícula,
% ejes y marcas en sus propias capas (y area style pone los ejes encima), y
% dentro de un axis los \draw, \fill y \node no admiten capa (van con los
% datos de main); por eso todo lo que se dibuja dentro del axis es un \addplot.
% Uso: \begin{axis}[…, dicapas] (al final de las opciones);
%   barras y áreas: \addplot[on layer=dicbarras, …];
%   sombreados:     \dicsombra[estilo]{x1}{y1}{x2}{y2} (rectángulo, en dicfondo,
%                   bajo las barras y sobre la cuadrícula);
%   cero y umbrales: \dicraya[estilo]{(x1,y1) (x2,y2)} (en axis lines);
%   líneas: \addplot normal (main); puntos: \addplot[on layer=dicpuntos, …]
%   (las marcas de cualquier serie van solas a dicpuntos);
%   rótulos: \dicrotulo[opciones del nodo]{(x,y)}{texto} (en axis descriptions),
%   siempre sin fondo (nada de fill=white: si chocan con una línea, se desplazan;
%   CRITERIOS §7, regla 4).
% Los \capa{n} del vídeo no cambian: siguen envolviendo instrucciones completas.
\pgfdeclarelayer{dicfondo}
\pgfdeclarelayer{dicbarras}
\pgfdeclarelayer{dicpuntos}
\pgfplotsset{
  % (el segundo argumento reparte cuadrícula, ejes y rótulos como el estándar)
  layers/dic/.define layer set={axis background, axis grid, dicfondo, dicbarras,
    axis ticks, axis lines, axis tick labels, pre main, main, dicpuntos,
    axis descriptions, axis foreground}{/pgfplots/layers/standard},
  % leyenda sin caja blanca: toma el fondo del gráfico, también si el gráfico va
  % dentro de un cuadro gris (CRITERIOS §7, regla 4, 26/10/2026)
  dicapas/.style={axis on top=false, set layers=dic, mark layer=dicpuntos,
    legend style={fill=none}},
  dicsuelto/.style={forget plot, stack plots=false, sharp plot, mark=none},
}
\newcommand{\dicraya}[2][]{%
  \addplot[on layer=axis lines, dicsuelto, diccero, #1] coordinates {#2};}
\newcommand{\dicsombra}[5][]{%
  \addplot[on layer=dicfondo, dicsuelto, draw=none, #1]
    coordinates {(#2,#3) (#4,#3) (#4,#5) (#2,#5)};}
\newcommand{\dicrotulo}[3][]{%
  \addplot[on layer=axis descriptions, dicsuelto, draw=none] coordinates {#2}
    node[pos=0, #1] {#3};}

%   \begin{document}
%   \begin{tikzpicture}[tcee]
%     \ejes{Q}{P}{6}{5}
%     \draw[curva1] (0.5,4.5) -- (5,0.8) node[etiqueta,right] {$D$};
%     \capa{2}{\draw[curva2] ...;}   % en el vídeo aparece en el paso 2
%   \end{tikzpicture}
%   \end{document}
%
% Paleta: la de la web (morado, dorado, salvia) más tres colores de apoyo
% distinguibles en blanco y negro por el trazo.
% ==========================================================================
\usepackage[T1]{fontenc}
\usepackage{newpxtext,newpxmath}
\usepackage{xcolor}
\usepackage{tikz,pgfplots}
\pgfplotsset{compat=1.18}
\usetikzlibrary{arrows.meta,calc,positioning,patterns,decorations.pathreplacing,intersections,shapes.misc,backgrounds}

\definecolor{tceemorado}{HTML}{5F2987}
\definecolor{tceedorado}{HTML}{B8860B}
\definecolor{tceeverde}{HTML}{3C7D3A}
\definecolor{tceeazul}{HTML}{1F5F99}
\definecolor{tceerojo}{HTML}{B22222}
\definecolor{tceenaranja}{HTML}{E07B28}
\definecolor{tceegris}{HTML}{767171}
\definecolor{tceesalvia}{HTML}{E2EFD9}

% Capas para el vídeo: \capa{n}{…} dibuja su contenido solo a partir del paso n.
% En el PDF, la web y Word se ve todo; generar-video.py compila el gráfico una
% vez por paso con \def\tceepaso{k} para que cada elemento aparezca al narrarlo.
% Vectores con flecha encima, como en el texto (tcee.sty): \vect{x}, \vect{p}
\providecommand{\vect}[1]{\vec{#1}}
\providecommand{\tceepaso}{1000}
\newcommand{\capa}[2]{\ifnum#1>\tceepaso\relax\else#2\fi}
% Gráficos por pasos del mismo tamaño (propuesta del 3.A.29, 11/11/2026): con
% \def\tceefantasma{1} antes del \input (p. ej., x-paso1.tex:
% «\def\tceepaso{3}\def\tceefantasma{1}\input{x.tex}»), las capas posteriores
% al paso se dibujan invisibles y ocupan su sitio. Sin \tceefantasma (el vídeo),
% no cambia nada.
\ifdefined\tceefantasma
  \renewcommand{\capa}[2]{\ifnum#1>\tceepaso\relax\begin{scope}[opacity=0]#2\end{scope}\else#2\fi}
\fi

\tikzset{
  tcee/.style={line cap=round,line join=round,font=\small,>={Stealth[length=2.2mm]}},
  eje/.style={->,thick,black},
  curva1/.style={very thick,tceemorado},
  curva2/.style={very thick,tceeazul},
  curva3/.style={very thick,tceeverde},
  curva4/.style={very thick,tceedorado},
  curvanueva/.style={very thick,tceerojo},
  desplazada/.style={very thick,dashed},
  guia/.style={thin,dashed,tceegris},
  punto/.style={circle,fill=black,inner sep=1.3pt},
  etiqueta/.style={font=\small,inner sep=2pt},
  flecha/.style={->,thick,tceerojo},
  area/.style={fill=tceemorado,fill opacity=0.18,draw=none},
  area2/.style={fill=tceedorado,fill opacity=0.22,draw=none},
}

% \ejes{etiqueta x}{etiqueta y}{largo x}{alto y}
\newcommand{\ejes}[4]{%
  \draw[eje] (0,0) -- (#3,0) node[below left=2pt and 0pt] {#1};
  \draw[eje] (0,0) -- (0,#4) node[left] {#2};
  \node[below left] at (0,0) {$0$};}
% \proyeccion{(x,y)}{etq x}{etq y}: líneas guía a los ejes con etiquetas
\newcommand{\proyeccion}[3]{%
  \draw[guia] let \p1=#1 in (\x1,0) node[below] {#2} |- (0,\y1) node[left] {#3};}

% --------------------------------------------------------------------------
% Paleta «dictamen» (dic…): gráficos de datos de la Guía del dictamen (1D01) y
% de cualquier gráfico de coyuntura que quiera los mismos colores. Un concepto,
% un color en todos los gráficos (CRITERIOS §7, «tres reglas», 24/10/2026);
% tabla completa en fuentes/1D01/_trabajo/paleta-1D01.md. No cambia nada de
% lo anterior: son colores y estilos nuevos.
% Actividad: el PIB (y sus variantes: PIB por persona, PIB mundial), morado.
\definecolor{dicpib}{HTML}{5F2987}
% Demanda interna, gama cálida: consumo (naranjas: privado oscuro, público
% claro), inversión (verde) y existencias (gris con trama, aparte).
\definecolor{dicconspriv}{HTML}{C76A1F}
\definecolor{dicconspub}{HTML}{F5C690}
\definecolor{dicfbcf}{HTML}{2A6B3C}
\definecolor{dicexistencias}{HTML}{B4B4B4}
% Demanda externa y comercio, gama fría (azul claro, que se separa del morado)
\definecolor{dicexterna}{HTML}{6AA6DA}
% Zonas: zona del euro (azul marino), España (rojo), resto de países (grises)
\definecolor{dicze}{HTML}{17406D}
\definecolor{dicespana}{HTML}{C8382E}
\definecolor{dicotro}{HTML}{8A8A8A}
\definecolor{dicotroclaro}{HTML}{B5B5B5}
\definecolor{dicotromasclaro}{HTML}{DADADA}
% Precios: inflación general (negro, continuo) y subyacente (negro,
% discontinuo); componentes: energía (bermellón), alimentos (amarillo) y los
% de la subyacente en azules (bienes industriales claro, servicios oscuro).
\definecolor{dicinflacion}{HTML}{1A1A1A}
\definecolor{dicenergia}{HTML}{D55E00}
\definecolor{dicalimentos}{HTML}{F2D45C}
\definecolor{dicbine}{HTML}{7FB2E0}
\definecolor{dicservicios}{HTML}{1F4E89}
% Trabajo: empleo (verde azulado), horas (el mismo, claro), tasa de actividad
% (verde azulado oscuro), paro (teja); productividad (dorado de la web).
\definecolor{dicempleo}{HTML}{1B7F6B}
\definecolor{dichoras}{HTML}{86CBB8}
\definecolor{dicactividad}{HTML}{0F5247}
\definecolor{dicparo}{HTML}{A23B2A}
\definecolor{dicproductividad}{HTML}{B8860B}
% Hogares: renta disponible (pizarra) y tasa de ahorro (negro); el consumo es
% dicconspriv.
\definecolor{dicrenta}{HTML}{4A5A6A}
\definecolor{dicahorro}{HTML}{1A1A1A}
% Política monetaria: tipos del BCE en azules (depósito, el de la zona del
% euro; operaciones principales, medio; facilidad marginal, claro); balance:
% préstamos a los bancos (dorado claro), resto del activo (gris claro).
\definecolor{dicdeposito}{HTML}{17406D}
\definecolor{dicopf}{HTML}{4F86C6}
\definecolor{dicfmc}{HTML}{9EC5EA}
\definecolor{dicprestamos}{HTML}{D9A520}
\definecolor{dicresto}{HTML}{CFCFCF}
% Ciclo: series adelantadas, coincidentes (el PIB) y retrasadas
\definecolor{dicadelantada}{HTML}{D55E00}
\definecolor{dicretrasada}{HTML}{6AA6DA}
% Referencias: rejilla, previsión, sombreados
\definecolor{dicreferencia}{HTML}{767171}

\tikzset{
  % barra de un componente: \addplot[dicbarra=dicfbcf]
  dicbarra/.style={fill=#1, draw=#1!65!black, line width=0.3pt},
  % existencias: gris con trama, para que se distinga también en gris
  dicbarratrama/.style={fill=#1, draw=#1!65!black, line width=0.3pt,
    postaction={pattern=north east lines, pattern color=white}},
  % línea de una serie; con halo blanco cuando pasa sobre barras
  diclinea/.style={very thick, #1},
  diclineahalo/.style={very thick, #1, preaction={draw=white, line width=3.4pt, solid}},
  % puntos de una serie (capa aparte, después de la línea)
  dicpuntos/.style={only marks, mark=*, mark size=1.8pt,
    mark options={fill=white, draw=#1, line width=0.9pt, solid}},
  % línea del cero o de un umbral, después de las barras
  diccero/.style={thin, black},
  dicumbral/.style={thin, dashed, dicreferencia},
}
% Capas de los gráficos de datos (CRITERIOS §7, «tres reglas»), de abajo
% arriba: cuadrícula → sombreados (dicfondo) → barras y áreas (dicbarras) → ejes, cero y umbrales
% (axis lines) → líneas (main) → puntos (dicpuntos) → rótulos (axis
% descriptions). No basta el orden del código: pgfplots dibuja cuadrícula,
% ejes y marcas en sus propias capas (y area style pone los ejes encima), y
% dentro de un axis los \draw, \fill y \node no admiten capa (van con los
% datos de main); por eso todo lo que se dibuja dentro del axis es un \addplot.
% Uso: \begin{axis}[…, dicapas] (al final de las opciones);
%   barras y áreas: \addplot[on layer=dicbarras, …];
%   sombreados:     \dicsombra[estilo]{x1}{y1}{x2}{y2} (rectángulo, en dicfondo,
%                   bajo las barras y sobre la cuadrícula);
%   cero y umbrales: \dicraya[estilo]{(x1,y1) (x2,y2)} (en axis lines);
%   líneas: \addplot normal (main); puntos: \addplot[on layer=dicpuntos, …]
%   (las marcas de cualquier serie van solas a dicpuntos);
%   rótulos: \dicrotulo[opciones del nodo]{(x,y)}{texto} (en axis descriptions),
%   siempre sin fondo (nada de fill=white: si chocan con una línea, se desplazan;
%   CRITERIOS §7, regla 4).
% Los \capa{n} del vídeo no cambian: siguen envolviendo instrucciones completas.
\pgfdeclarelayer{dicfondo}
\pgfdeclarelayer{dicbarras}
\pgfdeclarelayer{dicpuntos}
\pgfplotsset{
  % (el segundo argumento reparte cuadrícula, ejes y rótulos como el estándar)
  layers/dic/.define layer set={axis background, axis grid, dicfondo, dicbarras,
    axis ticks, axis lines, axis tick labels, pre main, main, dicpuntos,
    axis descriptions, axis foreground}{/pgfplots/layers/standard},
  % leyenda sin caja blanca: toma el fondo del gráfico, también si el gráfico va
  % dentro de un cuadro gris (CRITERIOS §7, regla 4, 26/10/2026)
  dicapas/.style={axis on top=false, set layers=dic, mark layer=dicpuntos,
    legend style={fill=none}},
  dicsuelto/.style={forget plot, stack plots=false, sharp plot, mark=none},
}
\newcommand{\dicraya}[2][]{%
  \addplot[on layer=axis lines, dicsuelto, diccero, #1] coordinates {#2};}
\newcommand{\dicsombra}[5][]{%
  \addplot[on layer=dicfondo, dicsuelto, draw=none, #1]
    coordinates {(#2,#3) (#4,#3) (#4,#5) (#2,#5)};}
\newcommand{\dicrotulo}[3][]{%
  \addplot[on layer=axis descriptions, dicsuelto, draw=none] coordinates {#2}
    node[pos=0, #1] {#3};}

\pgfplotsset{
  supu/.style={surf,z buffer=sort,shader=faceted,fill opacity=0.40,draw=tceemorado!35,line width=0.12pt,
    colormap={tceeu}{rgb255=(226,239,217) rgb255=(167,137,189)}},
  supg/.style={surf,z buffer=sort,shader=faceted,fill opacity=0.30,draw=tceeverde!35,line width=0.12pt,
    colormap={tceeg}{rgb255=(240,246,236) rgb255=(140,190,138)}},
  tresd/.style={x={(1.0cm,-0.275cm)},y={(0.575cm,0.467cm)},z={(0cm,1.0cm)},axis lines=none,
    xmin=0,xmax=4,ymin=0,ymax=4,zmin=0,zmax=4.4,clip=false,anchor=origin},
  % en el dual, el eje vertical (gasto) va a escala 0,6 para que el plano quepa
  tresdg/.style={tresd,z={(0cm,0.6cm)},zmax=7.6},
}
\begin{document}
\begin{tikzpicture}[tcee,
    cajap/.style={draw=tceemorado,very thick,fill=tceemorado!8,rounded corners=3pt,align=center,
                  minimum width=4.3cm,minimum height=1.05cm,inner sep=4pt},
    cajad/.style={cajap,draw=tceeverde,fill=tceeverde!8},
    cajaa/.style={draw=tceegris,thick,rounded corners=3pt,align=left,text width=5.15cm,
                  font=\footnotesize,inner sep=4pt,anchor=north},
    une/.style={->,thick},
    rel/.style={->,thick,tceegris},
    nota/.style={font=\footnotesize,inner sep=2pt},
    notag/.style={nota,text=tceegris},
    prop/.style={font=\footnotesize,align=left,inner sep=3pt,text width=5.5cm}]

  % ======================= cabeceras =======================
  \capa{1}{\node[text=tceemorado,align=center] at (4.4,0.75) {\textbf{Problema primal}\\ \footnotesize maximizar la utilidad con la renta $\overline{W}$ dada};}
  \capa{5}{\node[text=tceeverde,align=center] at (13.6,0.75) {\textbf{Problema dual}\\ \footnotesize minimizar el gasto con la utilidad $\overline{U}$ exigida};}

  % ======================= panel primal (3D) =======================
  \begin{axis}[tresd,at={(1.55cm,-5.1cm)}]
    % 1: ejes
    \capa{1}{\draw[eje] (axis cs:0,0,0) -- (axis cs:4.2,0,0) node[below] {$x_1$};}
    \capa{1}{\draw[eje] (axis cs:0,0,0) -- (axis cs:0,4.0,0) coordinate (finx2p);}
    \capa{1}{\draw[eje] (axis cs:0,0,0) -- (axis cs:0,0,4.3) node[above] {$U(\vect{x})$};}
    % 2: conjunto presupuestario y curva de indiferencia alcanzada (proyección en el suelo)
    \capa{2}{\addplot3[fill=tceedorado,fill opacity=0.25,draw=tceedorado,thick] coordinates {(0,0,0) (3.5,0,0) (0,3.5,0) (0,0,0)};}
    \capa{2}{\addplot3[tceemorado,thick,dashed,domain=1.0:3.06,samples=40,samples y=0] (x,{3.0625/x},0);}
    % 3: superficie de utilidad
    \capa{3}{\addplot3[supu,domain=0:3.6,y domain=0:3.6,samples=13,samples y=13] {sqrt(x*y)};}
    \capa{1}{\node[right] at (finx2p) {$x_2$};}
    % 4: restricción de gasto: plano vertical sobre la recta de balance
    \capa{4}{\addplot3[fill=tceedorado,fill opacity=0.22,draw=tceedorado,thin] coordinates {(3.5,0,0) (0,3.5,0) (0,3.5,2.35) (3.5,0,2.35) (3.5,0,0)};}
    % 5: la utilidad recorrida sobre la recta de balance y su máximo
    \capa{5}{\addplot3[tceemorado,line width=1.6pt,domain=0:3.5,samples=50,samples y=0] (x,{3.5-x},{sqrt(x*(3.5-x))});}
    \capa{5}{\draw[guia] (axis cs:1.75,1.75,1.75) -- (axis cs:1.75,1.75,0);}
    \capa{5}{\draw[guia] (axis cs:1.75,1.75,1.75) -- (axis cs:0,0,1.75);}
    \capa{5}{\node[punto,fill=tceemorado] at (axis cs:1.75,1.75,1.75) {};}
    \capa{5}{\node[punto,fill=tceemorado] at (axis cs:1.75,1.75,0) {};}
    \capa{5}{\node[nota,left,text=tceemorado] at (axis cs:0,0,1.75) {$U^*=V(\vect{p},\overline{W})$};}
    \capa{5}{\node[nota,left,text=tceemorado,inner sep=3pt] at (axis cs:1.75,1.75,0) {$\vect{x}^*$};}
    \capa{2}{\node[nota,text=tceedorado!80!black,anchor=north east] at (axis cs:2.1,0.9,0) {$\vect{p}\cdot\vect{x}\leq\overline{W}$};}
  \end{axis}

  % ======================= panel dual (3D) =======================
  \begin{axis}[tresdg,at={(10.75cm,-5.1cm)}]
    % 6: ejes
    \capa{6}{\draw[eje] (axis cs:0,0,0) -- (axis cs:4.2,0,0) node[below] {$x_1$};}
    \capa{6}{\draw[eje] (axis cs:0,0,0) -- (axis cs:0,4.0,0) coordinate (finx2d);}
    \capa{6}{\draw[eje] (axis cs:0,0,0) -- (axis cs:0,0,7.4) node[above] {$\vect{p}\cdot\vect{h}$};}
    % 7: curva de indiferencia exigida en el suelo y recta isogasto tangente
    \capa{7}{\addplot3[tceeverde,very thick,domain=1.0:3.06,samples=40,samples y=0] (x,{3.0625/x},0);}
    \capa{7}{\addplot3[tceeverde,thick,dashed,domain=0.6:2.9,samples=2,samples y=0] (x,{3.5-x},0);}
    % 8: plano del gasto
    \capa{8}{\addplot3[supg,domain=0:3.6,y domain=0:3.6,samples=9,samples y=9] {x+y};}
    \capa{6}{\node[right] at (finx2d) {$x_2$};}
    % 9: restricción de utilidad: «cortina» vertical sobre la curva de indiferencia
    % (camino de ida por el suelo y de vuelta por el plano, relleno)
    \capa{9}{\addplot3[fill=tceedorado,fill opacity=0.22,draw=none,domain=0:2,samples=121,samples y=0]
      ({x<1 ? 1.0+2.06*x : 1.0+2.06*(2-x)},{3.0625/(x<1 ? 1.0+2.06*x : 1.0+2.06*(2-x))},
       {x<1 ? 0 : (1.0+2.06*(2-x))+3.0625/(1.0+2.06*(2-x))});}
    % 10: el gasto recorrido sobre la curva de indiferencia y su mínimo
    \capa{10}{\addplot3[tceeverde,line width=1.6pt,domain=1.0:3.06,samples=50,samples y=0] (x,{3.0625/x},{x+3.0625/x});}
    \capa{10}{\draw[guia] (axis cs:1.75,1.75,3.5) -- (axis cs:1.75,1.75,0);}
    \capa{10}{\draw[guia] (axis cs:1.75,1.75,3.5) -- (axis cs:0,0,3.5);}
    \capa{10}{\node[punto,fill=tceeverde] at (axis cs:1.75,1.75,3.5) {};}
    \capa{10}{\node[punto,fill=tceeverde] at (axis cs:1.75,1.75,0) {};}
    \capa{10}{\node[nota,left,text=tceeverde] at (axis cs:0,0,3.5) {$E^*=E(\vect{p},\overline{U})$};}
    \capa{10}{\node[nota,left,text=tceeverde,inner sep=3pt] at (axis cs:1.75,1.75,0) {$\vect{h}^*$};}
    \capa{7}{\node[nota,text=tceeverde,anchor=west] at (axis cs:3.1,1.0,0) {$U(\vect{h})=\overline{U}$};}
  \end{axis}

  % lectura de cada panel
  \capa{5}{\node[nota,align=center,text width=6.4cm] at (4.4,-7.4) {sobre la recta de balance, la utilidad alcanza\\ su \emph{máximo} en $\vect{x}^*$: tangencia con $\overline{U}=U^*$};}
  \capa{10}{\node[nota,align=center,text width=6.4cm] at (13.6,-7.4) {sobre la curva de indiferencia $\overline{U}$, el gasto alcanza\\ su \emph{mínimo} en $\vect{h}^*$: tangencia con la isogasto $E^*$};}

  % ======================= esquema ampliado =======================
  % columnas: primal en x = 4,4; dual en x = 13,6
  \capa{11}{
    \node[cajap] (P) at (4.4,-8.65) {$\begin{array}{@{}r@{\quad}l@{}}\displaystyle\max_{\vect{x}} & U(\vect{x})\\[1pt] \text{s.a.} & \vect{p}\cdot\vect{x}\leq\overline{W},\ \vect{x}\geq\vect{0}\end{array}$};
    \node[nota,text=tceemorado,anchor=north] at (P.south) {$\mathcal{L}^{\mathrm{P}}=U(\vect{x})+\lambda(\overline{W}-\vect{p}\cdot\vect{x})$};
  }
  \capa{12}{
    \node[cajad] (D) at (13.6,-8.65) {$\begin{array}{@{}r@{\quad}l@{}}\displaystyle\min_{\vect{h}} & \vect{p}\cdot\vect{h}\\[1pt] \text{s.a.} & U(\vect{h})\geq\overline{U},\ \vect{h}\geq\vect{0}\end{array}$};
    \node[nota,text=tceeverde,anchor=north] at (D.south) {$\mathcal{L}^{\mathrm{D}}=\vect{p}\cdot\vect{h}+\mu\big(\overline{U}-U(\vect{h})\big)$};
  }
  % teorema básico
  \capa{13}{
    \draw[<->,very thick] (P.east) -- node[above,align=center,nota] {\textbf{teorema básico de la dualidad}\\ $\overline{W}=E^*,\ \overline{U}=U^*\ \Rightarrow\ \vect{x}^*=\vect{h}^*$} node[below,align=center,nota] {misma tangencia: $\dfrac{U_1}{U_2}=\dfrac{p_1}{p_2}$} (D.west);
  }
  % demandas
  \capa{14}{
    \node[cajap] (X) at (4.4,-11.6) {\textbf{Demanda marshalliana}\\ $\vect{x}^*=\vect{x}(\vect{p},\overline{W})$};
    \draw[une] ([yshift=-0.42cm]P.south) -- node[right,nota] {resolviendo} (X.north);
  }
  \capa{15}{
    \node[cajad] (H) at (13.6,-11.6) {\textbf{Demanda hicksiana}\\ $\vect{h}^*=\vect{h}(\vect{p},\overline{U})$};
    \draw[une] ([yshift=-0.42cm]D.south) -- node[left,nota] {resolviendo} (H.north);
  }
  % funciones valor
  \capa{16}{
    \node[cajap] (V) at (4.4,-14.9) {\textbf{Función indirecta de utilidad}\\ $V(\vect{p},\overline{W})=U\big(\vect{x}(\vect{p},\overline{W})\big)$};
    \draw[une] ([xshift=0.6cm]X.south) -- node[right,nota,align=left] {función\\ valor} ([xshift=0.6cm]V.north);
  }
  \capa{17}{
    \node[cajad] (E) at (13.6,-14.9) {\textbf{Función de gasto}\\ $E(\vect{p},\overline{U})=\vect{p}\cdot\vect{h}(\vect{p},\overline{U})$};
    \draw[une] ([xshift=-0.6cm]H.south) -- node[left,nota,align=right] {función\\ valor} ([xshift=-0.6cm]E.north);
  }
  % Roy y Shephard (por fuera)
  \capa{18}{
    \draw[une,tceemorado] ([xshift=-1.3cm]V.north) -- ([xshift=-1.3cm]X.south);
    \node[nota,text=tceemorado,align=right,anchor=east] at (2.95,-13.25) {Identidad\\ de \textsc{Roy}\\[2pt] $x_i=-\dfrac{\partial V/\partial p_i}{\partial V/\partial\overline{W}}$};
  }
  \capa{19}{
    \draw[une,tceeverde] ([xshift=1.3cm]E.north) -- ([xshift=1.3cm]H.south);
    \node[nota,text=tceeverde,align=left,anchor=west] at (15.05,-13.25) {Lema de\\ \textsc{Shephard}\\[2pt] $h_i=\dfrac{\partial E}{\partial p_i}$};
  }
  % Slutsky entre las demandas
  \capa{20}{
    \draw[<->,thick] (X.east) -- (H.west);
    \node[nota,above] at (9,-11.55) {$\dfrac{\partial x_i}{\partial p_j}=\underbrace{\dfrac{\partial h_i}{\partial p_j}}_{\text{ES}}\underbrace{-\,x_j\dfrac{\partial x_i}{\partial\overline{W}}}_{\text{ER}}$};
    \node[notag,below] at (9,-11.65) {ecuación de \textsc{Slutsky}};
  }
  % funciones valor inversas
  \capa{21}{
    \draw[rel,black] ([yshift=2.5mm]E.west) -- node[above,nota] {$V\big(\vect{p},E(\vect{p},\overline{U})\big)\equiv\overline{U}$} ([yshift=2.5mm]V.east);
    \draw[rel,black] ([yshift=-2.5mm]V.east) -- node[below,nota] {$E\big(\vect{p},V(\vect{p},\overline{W})\big)\equiv\overline{W}$} ([yshift=-2.5mm]E.west);
  }
  % identidades entre demandas (diagonales)
  \capa{22}{
    \draw[rel,dashed] (V.north east) -- node[pos=0.93,sloped,notag,anchor=north east,inner sep=1pt,font=\scriptsize,align=center] {$\vect{x}(\vect{p},\overline{W})\equiv{}$\\ $\vect{h}\big(\vect{p},V(\vect{p},\overline{W})\big)$} (H.south west);
    \draw[rel,dashed,<-] (X.south east) -- node[pos=0.07,sloped,notag,anchor=north west,inner sep=1pt,font=\scriptsize,align=center] {$\vect{h}(\vect{p},\overline{U})\equiv{}$\\ $\vect{x}\big(\vect{p},E(\vect{p},\overline{U})\big)$} (E.north west);
  }
  % propiedades y multiplicadores
  \capa{23}{
    \node[prop,text=tceemorado,anchor=north west] at (1.25,-15.55) {%
      $\vect{x}$ y $V$: grado 0 en $(\vect{p},\overline{W})$;\\
      \textsc{Walras}: $\vect{p}\cdot\vect{x}(\vect{p},\overline{W})=\overline{W}$;\\
      $V$: creciente en $\overline{W}$, no creciente en $\vect{p}$ y cuasiconvexa;\\
      $\lambda^*=\partial V/\partial\overline{W}$: utilidad marginal de la renta.};
  }
  \capa{24}{
    \node[prop,text=tceeverde,anchor=north east] at (16.75,-15.55) {%
      $\vect{h}$: grado 0 en $\vect{p}$; $E$: grado 1 en $\vect{p}$;\\
      sin exceso de utilidad: $U\big(\vect{h}(\vect{p},\overline{U})\big)=\overline{U}$;\\
      $E$: creciente en $\overline{U}$, no decreciente y cóncava en $\vect{p}$;\\
      $\mu^*=\partial E/\partial\overline{U}=1/\lambda^*$: coste marginal de la utilidad.};
    \node[nota,align=center,anchor=north] at (9,-15.75) {matriz de \textsc{Slutsky}\\[2pt] $S=\Big[\dfrac{\partial h_i}{\partial p_j}\Big]=D^2_{\vect{p}}E$\\[2pt] simétrica, semidefinida\\ negativa y con $S\vect{p}=\vect{0}$};
  }
  % aplicaciones (bloque III del tema)
  \capa{25}{
    \draw[tceegris,thin] (0.1,-18.45) -- (17.9,-18.45);
    \node[notag,anchor=south west] at (0.1,-18.45) {\itshape Para qué sirve (bloque III)};
    \node[cajaa] at (3.0,-18.65) {\textbf{Integrabilidad.} El esquema al revés: de $\vect{x}$ a $E$ (\textsc{Shephard} y dualidad entre demandas: un sistema de ecuaciones en derivadas parciales) y de $E$ a las preferencias. Exige \textsc{Walras} y $S$ simétrica y semidefinida negativa.};
    \node[cajaa] at (9.0,-18.65) {\textbf{Ecuación de \textsc{Slutsky}.} Derivar $h_i(\vect{p},\overline{U})\equiv x_i\big(\vect{p},E(\vect{p},\overline{U})\big)$ respecto a $p_j$ separa el efecto sustitución ($\leq0$ si $j=i$) del efecto renta: bienes normales, inferiores y \textsc{Giffen}; sustitutivos y complementarios.};
    \node[cajaa] at (15.0,-18.65) {\textbf{Medidas de bienestar.} $E$ traduce la utilidad a euros:\\ $\mathrm{VC}=\overline{W}-E(\vect{p}^1,\overline{U}^0)$,\\ $\mathrm{VE}=E(\vect{p}^0,\overline{U}^1)-\overline{W}$,\\ áreas bajo las hicksianas; índice de \textsc{Konüs}, $E(\vect{p}^1,\overline{U})/E(\vect{p}^0,\overline{U})$.};
  }
\end{tikzpicture}
\end{document}
